\clearpage
\VerdierPhysicalPageAnchor{0069}%
\VerdierSemanticAnchor{ch1-2-2-2-continuation}%
构成 $C$ 中的一个 $m$-重复形。

\VerdierSemanticAnchor{ch1-2-2-2-morphism}
现设 $f\colon Y\to Y'$ 为 $C$ 中 $\varphi$-可求和 $n$-重复形的态射。令

\VerdierSemanticAnchor{ch1-2-2-2-eq6}
\begin{equation}
  \left(\int_\varphi f\right)^\tau
  =\bigoplus_{\varphi(\sigma)=\tau}f^\sigma\colon
  \left(\int_\varphi Y\right)^\tau
  \longrightarrow
  \left(\int_\varphi Y'\right)^\tau.
  \tag{2.2.2.6}
\end{equation}

\VerdierSemanticAnchor{ch1-2-2-2-functoriality}
显然，$\left(\int_\varphi f\right)^\tau$ 与 $\int_\varphi Y$ 和
$\int_\varphi Y'$ 的微分交换。因此，这样的一族态射定义了 $m$-重复形的态射
$\int_\varphi f\colon\int_\varphi Y\to\int_\varphi Y'$。还显然有
$\int_\varphi fg=\int_\varphi f\,\int_\varphi g$ 以及
$\int_\varphi\operatorname{id}_Y=\operatorname{id}_{\int_\varphi Y}$。
因此，$\int_\varphi$ 是从 $\varphi$-可求和 $n$-重复形范畴到 $m$-重复形
范畴的函子；这个函子是加法的。

\VerdierSemanticAnchor{ch1-2-2-3}
\noindent\textbf{2.2.3.}\quad
设 $\varphi\colon[n]\to[m]$ 为映射，$\tau\in\mathbb Z[m]$，并设
$i\in[m]$。令 $E_1$（分别 $E_2$）为集合

\VerdierSemanticAnchor{ch1-2-2-3-fibers}
\[
  \begin{aligned}
    E_1&=\{\sigma_1\in\mathbb Z[n]\mid\varphi(\sigma_1)=\tau\},\\
    \text{（分别 }E_2
      &=\{\sigma_2\in\mathbb Z[n]\mid
          \varphi(\sigma_2)=\tau+e_i\}\text{）。}
  \end{aligned}
\]
最后，设 $Y$ 为 $C$ 中的一个 $n$-重复形。对
$\sigma_1\in E_1$、$\sigma_2\in E_2$，令

\VerdierSemanticAnchor{ch1-2-2-3-eq1}
\begin{equation}
  u_{\sigma_2,\sigma_1}(Y)
  =\begin{cases}
    \varepsilon(\tau,i)\varepsilon(\sigma_1,j)d_Y^{\sigma_1,j},
      &\text{若 }\sigma_2-\sigma_1=e_j,\\
    0,&\text{否则},
  \end{cases}
  \tag{2.2.3.1}
\end{equation}
从而得到一个 $E_2\times E_1$ 型矩阵

\VerdierSemanticAnchor{ch1-2-2-3-eq2}
\begin{equation}
  M(\tau,i,Y)
  =\bigl(u_{\sigma_2,\sigma_1}(Y)\bigr)_{
    \sigma_2\in E_2,\,\sigma_1\in E_1}.
  \tag{2.2.3.2}
\end{equation}

\VerdierSemanticAnchor{ch1-2-2-3-matrix-identification}
当 $Y$ 为 $\varphi$-可求和时，这就是态射
$d_{\int_\varphi Y}^{\tau,i}$ 的矩阵。

\VerdierSemanticAnchor{ch1-2-2-3-extensions}
注意，一般而言，这个矩阵的每一行与每一列都只有有限多个非零项。因此，当
$C$ 中可数乘积与直和都可表示时，函子 $\int_\varphi$ 至少有两种先验可能的
扩张到任意复形的方法：
